\begin{abstract}

Fault-tolerant quantum computing relies on quantum error correction (QEC) schemes whose performance can be characterized by a threshold: 
a critical physical error rate below which logical errors can be suppressed by increasing the code distance. 
In this thesis, we investigate the threshold for Steane-type error correction circuits which perform syndrome extraction by logical encoded ancilla qubits in a fault-tolerant and single-shot setting.
We determine the logical error rate under a circuit-level noise model with faults in both gate operations and measurements.
Decoding is performed using an efficient maximum likelihood (ML) decoder designed for the surface code under bit- or phase-flip noise and compared to minimum weight perfect matching decoder (MWPM). 
We find that the threshold for Steane-type error correction takes a value of approximately $3\%$ for a single QEC cycle, 
which falls of to a threshold value of $2\%$ for repeated independent QEC cycles.
Compared to a conventional repeated syndrome extraction circuit, which shows a threshold at approximately $1\%$ under circuit level noise,
we are able to show an improvement of factor $2$, respectively $3$. 
We furthermore investigate the influence of the order of the two independent half-cycles of Steane-type error correction, 
finding a significant asymmetry for the single QEC cycle case which is suppressed for the repeated cycle case. 
We perform those simulations using both a MWPM and ML decoder to understand the influence of the decoder on the performance. 
We find that the ML decoder in consistently performing better, even under restriction that we reduce the underlying error model to a phase- or bit-flip model. 

A limitation of this work is the assumption of access to fault tolerant prepared logical ancilla states, which warrants further investigation.

\end{abstract}
